% !PRINTED 596
\weberpage{596}
If one sets $h = p$, this yields
\[ \tau\mkern1mu|\mkern1mu\chi_1^p \equiv \tau\ (\mathrm{mod}\ \mathfrak{p}^3), \]
and in the same way one concludes from this, for every exponent $\nu$,
\weq{17}{\tau\mkern1mu|\mkern1mu\chi_1^{p^\nu} \equiv \tau\ (\mathrm{mod}\ \mathfrak{p}^{\nu+2}).}
To prove this formula in general, one applies complete induction. One assumes (17) as proved and sets
\[ \tau\mkern1mu|\mkern1mu\chi_1^{p^\nu} = \tau + \pi^{\nu+2}\alpha; \]
from this, by repeated application,
\[
\begin{aligned}
\tau\mkern1mu|\mkern1mu\chi_1^{2p^\nu} &\equiv \tau + 2\pi^{\nu+2}\alpha \\
\tau\mkern1mu|\mkern1mu\chi_1^{3p^\nu} &\equiv \tau + 3\pi^{\nu+2}\alpha\ (\mathrm{mod}\ \mathfrak{p}^{\nu+3}) \\
&\cdots\cdots\cdots\cdots \\
\tau\mkern1mu|\mkern1mu\chi_1^{p^{\nu+1}} &\equiv \tau + p\pi^{\nu+2}\alpha,
\end{aligned}
\]
and thus the formula (17) for $\nu + 1$.

Now one can always take $\nu$ so great that, of the differences
\[ \tau - \tau_1,\ \tau - \tau_2,\ \tau - \tau_3,\ \ldots, \]
if $\tau$ is different from $\tau_1, \tau_2, \tau_3, \ldots$, none is divisible by $\mathfrak{p}^{\nu-2}$, and then it follows from (17) that $\chi_1^{p^\nu}$ must be the identical substitution. \gesperrt{According to this, the degree of every element $\chi_1$ is a power of $p$, and consequently, by the Cauchy-Sylow theorem (\S.\,29, I.), no prime different from $p$ can divide the degree of $X^{(1)}$, as was to be proved.}

If, then, $g$ is not divisible by $p$, then $X^{(1)}$ is certainly the unit group.

One can continue in this way and form a group $X^{(2)}$ of degree $g_2$ from the congruence
\[ \tau\mkern1mu|\mkern1mu\chi_2 \equiv \tau\ (\mathrm{mod}\ \mathfrak{p}^3), \]
$X^{(2)}$ is a divisor of $X^{(1)}$, hence $g_2$ is likewise a power of $p$, and so one must finally arrive at the unit group. Among the successive groups $X^{(1)}, X^{(2)}, \ldots$ several may under certain circumstances also be equal to one another\footnote{\gesperrt{Hilbert} calls $\Psi$ the decomposition group, $X$ the inertia group, $X^{(1)}$ the ramification group of the ideal $\mathfrak{p}$. In the memoir mentioned above it is further proved that $g_1$ is the highest power of $p$ dividing $g$, that $X^{(1)}$ is a normal divisor of $X$, and that the whole group $\Psi$ is metacyclic.}.
